\section{Problem Formulation}
\label{sec:setting}

Let $\mathcal D_d$ be the prompt distribution for domain $d$, $q_d(v\mid h)$ its teacher distribution, and $p_\theta(v\mid h)$ the shared student distribution over vocabulary $\mathcal V$. At each update, the student samples a response $y$ to prompt $x$. Its assigned teacher scores the same prefixes $h=(x,y_{<t})$. Let $\lambda_d\geq0$ be the target domain weights, with $\sum_d\lambda_d=1$.

For a per-position distillation loss $\ell_{d,t}$, a response-normalized objective is
\begin{equation}
\mathcal L(\theta)=\sum_d\lambda_d\,
\mathbb E_{x\sim\mathcal D_d,\,y\sim p_\theta(\cdot\mid x)}
\left[\frac{1}{|y|}\sum_{t=1}^{|y|}\ell_{d,t}(\theta)\right].
\label{eq:mopd}
\end{equation}
Training differentiates the losses on the sampled responses while holding prefixes, valid-token masks, and response weights fixed. Section~\ref{sec:normalization} compares the resulting batch gradient with those from other averaging rules.

At a fixed prefix, a full-vocabulary reference loss is reverse KL:
\begin{equation}
\ell_{\mathrm{full}}=\sum_{v\in\mathcal V}p_\theta(v)\log\frac{p_\theta(v)}{q_d(v)}.
\label{eq:full}
\end{equation}
A sampled-token surrogate, written using stop-gradient $\sg[\cdot]$, is
\begin{equation}
\ell_{\mathrm{PG}}(a)=-\sg[\log q_d(a)-\log p_\theta(a)]\log p_\theta(a),
\qquad a\sim p_\theta.
\label{eq:pg}
\end{equation}
At a fixed prefix, its expected gradient equals $\nabla_\theta\ell_{\mathrm{full}}$; Appendix~\ref{app:losses} gives the derivation. Vocabulary-restricted supervision instead evaluates a selected set of alternatives. We use \emph{supervision density} to refer to the number of vocabulary alternatives scored at each prefix, and \emph{update concentration} to describe how parameter changes are distributed across coordinates. These refer to different objects.
